\documentclass[11pt,a4paper]{article}
\usepackage[T1]{fontenc}
\usepackage[utf8]{inputenc}
\usepackage[italian]{babel}
\usepackage{amsmath,amssymb,enumitem,xcolor}
\usepackage[a4paper,margin=2.2cm]{geometry}
\newcounter{quiz}
\newenvironment{quiz}{\par\medskip\noindent\refstepcounter{quiz}\textbf{\thequiz)}\enspace\ignorespaces}{\par\bigskip}
\newcommand{\registraSoluzione}[2]{}
\begin{document}
\section*{Quiz Analisi 1 ÔÇö Paola Suria}
\input{quiz/quiz_pagine_001_010.tex}
\input{quiz/quiz_pagine_011_020.tex}
\input{quiz/quiz_pagine_021_030.tex}
\input{quiz/quiz_pagine_031_040.tex}
\input{quiz/quiz_pagine_041_050.tex}
\input{quiz/quiz_pagine_051_060.tex}
\input{quiz/quiz_pagine_061_070.tex}
\input{quiz/quiz_pagine_071_080.tex}
\input{quiz/quiz_pagine_081_090.tex}
\input{quiz/quiz_pagine_091_100.tex}
\input{quiz/quiz_pagine_101_110.tex}
\input{quiz/quiz_pagine_111_120.tex}
\input{quiz/quiz_pagine_121_130.tex}
\input{quiz/quiz_pagine_131_140.tex}
\input{quiz/quiz_pagine_141_150.tex}
\input{quiz/quiz_pagine_151_160.tex}
\input{quiz/quiz_pagine_161_170.tex}
\input{quiz/quiz_pagine_171_180.tex}
\input{quiz/quiz_pagine_181_190.tex}
\input{quiz/quiz_pagine_191_200.tex}
\input{quiz/quiz_pagine_201_210.tex}
\input{quiz/quiz_pagine_211_220.tex}
\input{quiz/quiz_pagine_221_230.tex}
\input{quiz/quiz_pagine_231_240.tex}
\input{quiz/quiz_pagine_241_250.tex}
\input{quiz/quiz_pagine_251_260.tex}
\input{quiz/quiz_pagine_261_270.tex}
\input{quiz/quiz_pagine_271_280.tex}
\input{quiz/quiz_pagine_281_290.tex}
\input{quiz/quiz_pagine_291_300.tex}
\input{quiz/quiz_pagine_301_310.tex}
\input{quiz/quiz_pagine_311_320.tex}
\input{quiz/quiz_pagine_321_330.tex}
\input{quiz/quiz_pagine_331_340.tex}
\input{quiz/quiz_pagine_341_350.tex}
\input{quiz/quiz_pagine_351_360.tex}
\input{quiz/quiz_pagine_361_370.tex}
\input{quiz/quiz_pagine_371_380.tex}
\input{quiz/quiz_pagine_381_390.tex}
\input{quiz/quiz_pagine_391_400.tex}
\input{quiz/quiz_pagine_401_410.tex}
\input{quiz/quiz_pagine_411_420.tex}
\input{quiz/quiz_pagine_421_430.tex}
\input{quiz/quiz_pagine_431_440.tex}
\input{quiz/quiz_pagine_441_450.tex}
\input{quiz/quiz_pagine_451_460.tex}
\input{quiz/quiz_pagine_461_470.tex}
\input{quiz/quiz_pagine_471_480.tex}
\input{quiz/quiz_pagine_481_490.tex}
\input{quiz/quiz_pagine_491_500.tex}
\input{quiz/quiz_pagine_501_510.tex}
\input{quiz/quiz_pagine_511_520.tex}
\input{quiz/quiz_pagine_521_521.tex}
\clearpage\section*{Soluzioni e spiegazioni per capitolo (AM1 Morandotti)}
\subsection*{Capitolo 1}
\subsubsection*{1.4 ÔÇö Funzioni}
\noindent\textbf{quiz\_pagine\_001\_010\_003}: B. Spiegazione non ancora disponibile.\par\medskip
\noindent\textbf{quiz\_pagine\_001\_010\_004}: B. Spiegazione non ancora disponibile.\par\medskip
\noindent\textbf{quiz\_pagine\_001\_010\_009}: E. Spiegazione non ancora disponibile.\par\medskip
\noindent\textbf{quiz\_pagine\_001\_010\_015}: A. Spiegazione non ancora disponibile.\par\medskip
\noindent\textbf{quiz\_pagine\_001\_010\_016}: E. Spiegazione non ancora disponibile.\par\medskip
\noindent\textbf{quiz\_pagine\_001\_010\_020}: E. Spiegazione non ancora disponibile.\par\medskip
\noindent\textbf{quiz\_pagine\_001\_010\_028}: E. Spiegazione non ancora disponibile.\par\medskip
\noindent\textbf{quiz\_pagine\_001\_010\_030}: B. Spiegazione non ancora disponibile.\par\medskip
\noindent\textbf{quiz\_pagine\_001\_010\_031}: A. Spiegazione non ancora disponibile.\par\medskip
\noindent\textbf{quiz\_pagine\_001\_010\_035}: D. Spiegazione non ancora disponibile.\par\medskip
\noindent\textbf{quiz\_pagine\_001\_010\_036}: E. Spiegazione non ancora disponibile.\par\medskip
\noindent\textbf{quiz\_pagine\_001\_010\_040}: B. Spiegazione non ancora disponibile.\par\medskip
\noindent\textbf{quiz\_pagine\_011\_020\_002}: B. Su $[0,+\infty)$ \textbackslash{}`e strettamente decrescente, dunque iniettiva.\par\medskip
\noindent\textbf{quiz\_pagine\_011\_020\_010}: D. $(f^{-1})\prime(2)=1/f\prime(0)=1/(2+3)=1/5$.\par\medskip
\noindent\textbf{quiz\_pagine\_011\_020\_011}: D. $e^{x^4}=1+x^4+o(x^4)$ e $\sqrt{1+x^4}=1+x^4/2+o(x^4)$; la differenza \textbackslash{}`e positiva.\par\medskip
\noindent\textbf{quiz\_pagine\_011\_020\_027}: E. $w=(5+i)(3+i)/10=(14+8i)/10$, parte reale $7/5$. Nel testo \textbackslash{}`e scritto $\operatorname{Re}(z)$ anzich\textbackslash{}'e $\operatorname{Re}(w)$.\par\medskip
\noindent\textbf{quiz\_pagine\_011\_020\_028}: E. $f\prime(x)=5x^4+3e^x>0$; $f(0)=1>0$ e $f(x)\to-\infty$ a sinistra.\par\medskip
\noindent\textbf{quiz\_pagine\_011\_020\_031}: B. $|w|=\sqrt{26}/\sqrt{10}=\sqrt{65}/5$.\par\medskip
\noindent\textbf{quiz\_pagine\_021\_030\_006}: E. Il radicando deve essere strettamente positivo, $1-x^2>0$, e $x+4>0$.\par\medskip
\noindent\textbf{quiz\_pagine\_021\_030\_007}: E. $49x^{48}+4>0$, quindi esiste un solo zero reale, positivo. Le radici non sono in generale equidistanti su una circonferenza.\par\medskip
\noindent\textbf{quiz\_pagine\_021\_030\_010}: D. $|z|z=\sqrt{13}(2+3i)$ non \textbackslash{}`e reale.\par\medskip
\noindent\textbf{quiz\_pagine\_021\_030\_011}: E. La base \textbackslash{}`e $e^{i\pi/3}$; elevando alla nona si ottiene $e^{i3\pi}=-1$.\par\medskip
\noindent\textbf{quiz\_pagine\_021\_030\_014}: A. $g\prime=-\sin(f)f\prime$: i tre zeri di $f$ e i due zeri distinti di $f\prime$ forniscono almeno cinque punti critici.\par\medskip
\noindent\textbf{quiz\_pagine\_021\_030\_016}: non verificata / nessuna alternativa corretta. I termini dominanti sono $\sqrt{x}$ al numeratore e $-4\sqrt{x}$ al denominatore: il limite risulterebbe $-1/4$. La lettura dell\textbackslash{}'espressione nel PDF richiede verifica.\par\medskip
\noindent\textbf{quiz\_pagine\_021\_030\_017}: C. $f(0)=3$, $f\prime(0)=6$, quindi $(f^{-1})\prime(3)=1/6$.\par\medskip
\noindent\textbf{quiz\_pagine\_021\_030\_018}: C. Vicino a $x=1$ il termine $|\cos(\pi x)|^3$ ha derivata nulla nel punto, mentre $2|x|$ ha derivata $2$.\par\medskip
\noindent\textbf{quiz\_pagine\_021\_030\_020}: D. Con $u=2-3x^2$, $du=-6x\,dx$, perci\textbackslash{}`o l\textbackslash{}'integrale \textbackslash{}`e $-u^6/36+C$.\par\medskip
\noindent\textbf{quiz\_pagine\_021\_030\_024}: D. Le radici hanno argomenti $-\pi/6$, $\pi/2$, $7\pi/6$: nessuna coppia \textbackslash{}`e coniugata.\par\medskip
\noindent\textbf{quiz\_pagine\_021\_030\_026}: E. La costante \textbackslash{}`e $\pi/2$ per $x>0$ e $-\pi/2$ per $x<0$, quindi non \textbackslash{}`e unica su $\mathbb R\setminus\{0\}$.\par\medskip
\noindent\textbf{quiz\_pagine\_021\_030\_030}: E. Poich\textbackslash{}'e $-1\le\cos x\le1$, si ha $1/2\le2^{\cos x}\le2$.\par\medskip
\noindent\textbf{quiz\_pagine\_021\_030\_033}: B. $\int_1^{+\infty}3x^{-3}dx=[-3/(2x^2)]_1^{+\infty}=3/2$.\par\medskip
\noindent\textbf{quiz\_pagine\_021\_030\_034}: D. L\textbackslash{}'integrale \textbackslash{}`e $-\log(\cos\pi/4)=\frac12\log2$; (d) vale il doppio.\par\medskip
\noindent\textbf{quiz\_pagine\_021\_030\_035}: A. $\sqrt[3]{1+3x}=1+x+o(x)$, quindi il rapporto tende a $1/7$.\par\medskip
\noindent\textbf{quiz\_pagine\_031\_040\_001}: E. Spiegazione non ancora disponibile.\par\medskip
\noindent\textbf{quiz\_pagine\_031\_040\_011}: D. Spiegazione non ancora disponibile.\par\medskip
\noindent\textbf{quiz\_pagine\_031\_040\_027}: A. Spiegazione non ancora disponibile.\par\medskip
\noindent\textbf{quiz\_pagine\_031\_040\_028}: A. Spiegazione non ancora disponibile.\par\medskip
\noindent\textbf{quiz\_pagine\_031\_040\_031}: A. Spiegazione non ancora disponibile.\par\medskip
\noindent\textbf{quiz\_pagine\_031\_040\_032}: B. Spiegazione non ancora disponibile.\par\medskip
\noindent\textbf{quiz\_pagine\_031\_040\_033}: B. Spiegazione non ancora disponibile.\par\medskip
\noindent\textbf{quiz\_pagine\_031\_040\_034}: B. Spiegazione non ancora disponibile.\par\medskip
\noindent\textbf{quiz\_pagine\_031\_040\_035}: A. Spiegazione non ancora disponibile.\par\medskip
\noindent\textbf{quiz\_pagine\_031\_040\_036}: A. Spiegazione non ancora disponibile.\par\medskip
\noindent\textbf{quiz\_pagine\_031\_040\_037}: B. Spiegazione non ancora disponibile.\par\medskip
\noindent\textbf{quiz\_pagine\_041\_050\_001}: B. Risposta ufficiale indicata nella tabella della decima simulazione.\par\medskip
\noindent\textbf{quiz\_pagine\_041\_050\_004}: E. Risposta ufficiale indicata nella tabella della decima simulazione.\par\medskip
\noindent\textbf{quiz\_pagine\_041\_050\_016}: A. Risposta ufficiale indicata nella tabella della decima simulazione.\par\medskip
\noindent\textbf{quiz\_pagine\_041\_050\_021}: non verificata / nessuna alternativa corretta. Soluzione non riportata nel PDF; da verificare.\par\medskip
\noindent\textbf{quiz\_pagine\_041\_050\_022}: non verificata / nessuna alternativa corretta. Soluzione non riportata nel PDF; da verificare.\par\medskip
\noindent\textbf{quiz\_pagine\_041\_050\_023}: non verificata / nessuna alternativa corretta. Soluzione non riportata nel PDF; da verificare.\par\medskip
\noindent\textbf{quiz\_pagine\_041\_050\_024}: non verificata / nessuna alternativa corretta. Soluzione non riportata nel PDF; da verificare.\par\medskip
\noindent\textbf{quiz\_pagine\_041\_050\_025}: non verificata / nessuna alternativa corretta. Soluzione non riportata nel PDF; da verificare.\par\medskip
\noindent\textbf{quiz\_pagine\_041\_050\_026}: non verificata / nessuna alternativa corretta. Soluzione non riportata nel PDF; da verificare.\par\medskip
\noindent\textbf{quiz\_pagine\_041\_050\_030}: non verificata / nessuna alternativa corretta. Formula verificata visivamente sulla pagina 8 del PDF.\par\medskip
\noindent\textbf{quiz\_pagine\_041\_050\_031}: non verificata / nessuna alternativa corretta. Soluzione non riportata nel PDF; da verificare.\par\medskip
\noindent\textbf{quiz\_pagine\_051\_060\_004}: C. Una funzione strettamente crescente è iniettiva e la sua inversa sull’immagine è strettamente crescente.\par\medskip
\noindent\textbf{quiz\_pagine\_051\_060\_008}: A. La composizione di due funzioni dispari è dispari.\par\medskip
\noindent\textbf{quiz\_pagine\_061\_070\_001}: E. Le prime quattro funzioni non sono pari; la risposta ufficiale è E.\par\medskip
\noindent\textbf{quiz\_pagine\_061\_070\_021}: B. Il dominio comune è x>1; su tale dominio vale la proprietà della radice del quoziente.\par\medskip
\noindent\textbf{quiz\_pagine\_071\_080\_020}: non verificata / nessuna alternativa corretta. Dato $f(0)=-2$, ogni minimo globale, se esiste, è necessariamente $\le-2$. Senza continuità l'esistenza del minimo non è garantita.\par\medskip
\noindent\textbf{quiz\_pagine\_091\_100\_001}: B. Ponendo $y=\log x+2$, si ottiene $x=e^{y-2}$; quindi $f^{-1}(x)=e^{x-2}$.\par\medskip
\noindent\textbf{quiz\_pagine\_091\_100\_018}: D. $y=\sqrt[7]{x+7}$ implica $y^7=x+7$, dunque $x=y^7-7$.\par\medskip
\noindent\textbf{quiz\_pagine\_101\_110\_001}: E. Ponendo $y=e^x+2$ si ottiene $x=\log(y-2)$, per $y>2$.\par\medskip
\noindent\textbf{quiz\_pagine\_101\_110\_018}: D. La composizione $g\circ f$ è $e^{4\log x}=x^4$, ma il suo dominio è $(0,+\infty)$, non $\mathbb R$.\par\medskip
\subsubsection*{1.5 ÔÇö Insiemi numerici}
\noindent\textbf{quiz\_pagine\_001\_010\_002}: A. Spiegazione non ancora disponibile.\par\medskip
\noindent\textbf{quiz\_pagine\_001\_010\_018}: D. Spiegazione non ancora disponibile.\par\medskip
\noindent\textbf{quiz\_pagine\_001\_010\_019}: E. Spiegazione non ancora disponibile.\par\medskip
\noindent\textbf{quiz\_pagine\_001\_010\_021}: C. Spiegazione non ancora disponibile.\par\medskip
\noindent\textbf{quiz\_pagine\_011\_020\_001}: C. $\cosh x>4$ equivale a $|x|>\operatorname{arcosh}4$: il dominio \textbackslash{}`e unione di due intervalli.\par\medskip
\noindent\textbf{quiz\_pagine\_011\_020\_003}: A. Per $z=x\in\mathbb R$, $(x-1)^2+1<9$, cio\textbackslash{}`e $1-\sqrt8<x<1+\sqrt8$.\par\medskip
\noindent\textbf{quiz\_pagine\_011\_020\_009}: B. Poich\textbackslash{}'e $f(0)>0$, vicino a zero $|f|=f$ e quindi $g\prime(0)=-4$.\par\medskip
\noindent\textbf{quiz\_pagine\_011\_020\_013}: D. $f\prime(x)=5x^4+\cosh x>0$, quindi $f$ \textbackslash{}`e strettamente crescente e il valore inferiore a $f(10)$ \textbackslash{}`e assunto prima di $10$.\par\medskip
\noindent\textbf{quiz\_pagine\_011\_020\_014}: A. $g\prime=3f^2f\prime$ si annulla in $0$, in $1$ e in un punto interno per Rolle.\par\medskip
\noindent\textbf{quiz\_pagine\_011\_020\_036}: E. $\log(1+x^2)\sim x^2$ e $x\sin x\sim x^2$; per $\alpha>2$ il limite \textbackslash{}`e $1$, per $\alpha=2$ \textbackslash{}`e $2$, altrimenti non \textbackslash{}`e $0$.\par\medskip
\noindent\textbf{quiz\_pagine\_021\_030\_002}: C. $f\prime$ ha segno di $-6x(x+1)-1$; \textbackslash{}`e positivo tra le radici $(-3\pm\sqrt3)/6$. Attenzione: l\textbackslash{}'intervallo attraversa $x=-1$? No, entrambe le radici sono maggiori di $-1$.\par\medskip
\noindent\textbf{quiz\_pagine\_021\_030\_005}: C. Ogni complesso non nullo ha esattamente quattro radici quarte distinte. La formula degli argomenti si intende modulo $2\pi$.\par\medskip
\noindent\textbf{quiz\_pagine\_021\_030\_013}: E. Moltiplicando per $1-i$ si ottiene $z=(1-i)^2/2=-i$.\par\medskip
\noindent\textbf{quiz\_pagine\_021\_030\_015}: C. Applicando Rolle a due punti con lo stesso valore, oppure usando un estremo interno quando $g$ non \textbackslash{}`e costante, si trova $g\prime(x_0)=0$.\par\medskip
\noindent\textbf{quiz\_pagine\_021\_030\_027}: B. Per Rolle applicato a $f-g$ esiste $c$ con $f\prime(c)=g\prime(c)$; basta prendere $c_1=c_2=c$.\par\medskip
\noindent\textbf{quiz\_pagine\_021\_030\_031}: C. Per definizione $[f(x)^3-f(0)^3]/(x-0)=f(x)^3/x$.\par\medskip
\noindent\textbf{quiz\_pagine\_021\_030\_032}: E. $(2-3i)^2=-5-12i$ e $i^7=-i$, quindi $(-i)(-5-12i)=-12+5i$. L\textbackslash{}'opzione corretta \textbackslash{}`e (e), non (a).\par\medskip
\noindent\textbf{quiz\_pagine\_031\_040\_012}: C. Spiegazione non ancora disponibile.\par\medskip
\noindent\textbf{quiz\_pagine\_031\_040\_014}: D. Spiegazione non ancora disponibile.\par\medskip
\noindent\textbf{quiz\_pagine\_031\_040\_015}: D. Spiegazione non ancora disponibile.\par\medskip
\noindent\textbf{quiz\_pagine\_031\_040\_016}: E. Spiegazione non ancora disponibile.\par\medskip
\noindent\textbf{quiz\_pagine\_031\_040\_017}: A. Spiegazione non ancora disponibile.\par\medskip
\noindent\textbf{quiz\_pagine\_031\_040\_018}: C. Spiegazione non ancora disponibile.\par\medskip
\noindent\textbf{quiz\_pagine\_031\_040\_029}: D. Spiegazione non ancora disponibile.\par\medskip
\noindent\textbf{quiz\_pagine\_031\_040\_030}: B. Spiegazione non ancora disponibile.\par\medskip
\noindent\textbf{quiz\_pagine\_041\_050\_005}: A. Risposta ufficiale indicata nella tabella della decima simulazione.\par\medskip
\noindent\textbf{quiz\_pagine\_061\_070\_004}: D. Poiché z coniugato(z)=|z|²=1, il luogo è la circonferenza unitaria.\par\medskip
\noindent\textbf{quiz\_pagine\_061\_070\_025}: A. La parte immaginaria di (a+ib)(3+i) è a+3b; si annulla per a=-3b.\par\medskip
\noindent\textbf{quiz\_pagine\_071\_080\_017}: E. Le radici cubiche di $-27$ sono $-3$ e $3/2\pm(3\sqrt3/2)i$; le ultime due sono coniugate.\par\medskip
\noindent\textbf{quiz\_pagine\_081\_090\_013}: A. $\overline z=-1-2i$, quindi $\overline z^{\,2}+3=4i$ e il suo modulo è $4$.\par\medskip
\noindent\textbf{quiz\_pagine\_081\_090\_014}: C. $\frac{1-2i}{3+4i}=\frac{(1-2i)(3-4i)}{25}=\frac{-5-10i}{25}=-\frac15-\frac25i$.\par\medskip
\noindent\textbf{quiz\_pagine\_081\_090\_017}: E. Fra le sei radici complesse, le due reali sono $z=\pm\sqrt[6]{11}$.\par\medskip
\noindent\textbf{quiz\_pagine\_091\_100\_004}: B. Il massimo è $5$; $1+1/(3n)>1$ e tende a $1$, che non appartiene ad $A$. Quindi non c'è minimo.\par\medskip
\noindent\textbf{quiz\_pagine\_091\_100\_005}: D. Per esempio $A=(0,1)$ e $B=[0,1]$ sono distinti, con $A\subset B$, ma hanno entrambi $\inf=0$ e $\sup=1$.\par\medskip
\noindent\textbf{quiz\_pagine\_091\_100\_026}: E. Ogni minorante di $B$ è minorante di $A$; quindi, per insiemi non vuoti e inferiormente limitati, $\inf(A)\geq\inf(B)$.\par\medskip
\noindent\textbf{quiz\_pagine\_101\_110\_013}: D. Le radici di $i$ sono $\pm(1+i)/\sqrt2$, cioè opposte e non coniugate. Entrambe hanno modulo $1$.\par\medskip
\subsection*{Capitolo 2}
\subsubsection*{2.1 ÔÇö Successioni reali e loro propriet├á}
\noindent\textbf{quiz\_pagine\_031\_040\_013}: C. Spiegazione non ancora disponibile.\par\medskip
\noindent\textbf{quiz\_pagine\_031\_040\_020}: E. Spiegazione non ancora disponibile.\par\medskip
\noindent\textbf{quiz\_pagine\_031\_040\_021}: A. Spiegazione non ancora disponibile.\par\medskip
\noindent\textbf{quiz\_pagine\_031\_040\_022}: E. Spiegazione non ancora disponibile.\par\medskip
\noindent\textbf{quiz\_pagine\_031\_040\_039}: A. Spiegazione non ancora disponibile.\par\medskip
\noindent\textbf{quiz\_pagine\_031\_040\_040}: B. Spiegazione non ancora disponibile.\par\medskip
\noindent\textbf{quiz\_pagine\_031\_040\_042}: B. Spiegazione non ancora disponibile.\par\medskip
\subsubsection*{2.2 ÔÇö Limiti di successioni}
\noindent\textbf{quiz\_pagine\_001\_010\_006}: A. Spiegazione non ancora disponibile.\par\medskip
\noindent\textbf{quiz\_pagine\_001\_010\_010}: E. Spiegazione non ancora disponibile.\par\medskip
\noindent\textbf{quiz\_pagine\_001\_010\_025}: B. Spiegazione non ancora disponibile.\par\medskip
\noindent\textbf{quiz\_pagine\_011\_020\_007}: A. $\cos(\pi n^2)=(-1)^n$, dunque $a_n=(-1)^n n$ oscilla senza limite.\par\medskip
\noindent\textbf{quiz\_pagine\_021\_030\_019}: A. Entrambi gli addendi sono non negativi, dunque ciascuno tende a zero.\par\medskip
\noindent\textbf{quiz\_pagine\_021\_030\_029}: B. $\sin(\pi n/2)$ assume ciclicamente $1,0,-1,0$, quindi la successione diverge oscillando e non ha limite.\par\medskip
\noindent\textbf{quiz\_pagine\_031\_040\_019}: A. Spiegazione non ancora disponibile.\par\medskip
\noindent\textbf{quiz\_pagine\_031\_040\_023}: E. Spiegazione non ancora disponibile.\par\medskip
\noindent\textbf{quiz\_pagine\_031\_040\_024}: D. Spiegazione non ancora disponibile.\par\medskip
\noindent\textbf{quiz\_pagine\_031\_040\_025}: C. Spiegazione non ancora disponibile.\par\medskip
\noindent\textbf{quiz\_pagine\_031\_040\_038}: B. Spiegazione non ancora disponibile.\par\medskip
\noindent\textbf{quiz\_pagine\_031\_040\_041}: B. Spiegazione non ancora disponibile.\par\medskip
\noindent\textbf{quiz\_pagine\_031\_040\_043}: A. Spiegazione non ancora disponibile.\par\medskip
\noindent\textbf{quiz\_pagine\_041\_050\_012}: C. Il rapporto tra radice cubica e termine lineare tende a zero; il seno è limitato.\par\medskip
\noindent\textbf{quiz\_pagine\_041\_050\_027}: non verificata / nessuna alternativa corretta. Soluzione non riportata nel PDF; da verificare.\par\medskip
\noindent\textbf{quiz\_pagine\_041\_050\_028}: non verificata / nessuna alternativa corretta. Soluzione non riportata nel PDF; da verificare.\par\medskip
\noindent\textbf{quiz\_pagine\_041\_050\_029}: non verificata / nessuna alternativa corretta. Soluzione non riportata nel PDF; da verificare.\par\medskip
\noindent\textbf{quiz\_pagine\_071\_080\_007}: A. Poiché $\sin(\pi/6)=1/2$ e l'esponente è negativo e tende a $-\infty$, la successione diverge a $+\infty$.\par\medskip
\noindent\textbf{quiz\_pagine\_081\_090\_006}: C. $b_n=c_n-a_n$ è la differenza fra una successione convergente e una che tende a $+\infty$; dunque $b_n\to-\infty$.\par\medskip
\noindent\textbf{quiz\_pagine\_081\_090\_016}: C. $\sin(5/(1+n))\sim5/n$ e $(\log n)/n\to0$, quindi il prodotto tende a $0$.\par\medskip
\noindent\textbf{quiz\_pagine\_101\_110\_004}: B. Il logaritmo vale $\log(1+\sqrt n)/n^2\to0$, quindi il limite è $e^0=1$.\par\medskip
\noindent\textbf{quiz\_pagine\_101\_110\_015}: D. Per il teorema del confronto, da $0\leq a_n\leq b_n\to0$ segue $a_n\to0$.\par\medskip
\subsection*{Capitolo 3}
\subsubsection*{3.1 ÔÇö Definizioni e prime propriet├á. Limiti}
\noindent\textbf{quiz\_pagine\_001\_010\_022}: D. Spiegazione non ancora disponibile.\par\medskip
\noindent\textbf{quiz\_pagine\_001\_010\_026}: C. Spiegazione non ancora disponibile.\par\medskip
\noindent\textbf{quiz\_pagine\_011\_020\_004}: D. $x^3\log(1+1/(4x^3))\to1/4$ per $\log(1+t)\sim t$.\par\medskip
\noindent\textbf{quiz\_pagine\_011\_020\_005}: E. $x[(1+2/x)^{1/4}-1]\to 1/2$.\par\medskip
\noindent\textbf{quiz\_pagine\_031\_040\_003}: C. Spiegazione non ancora disponibile.\par\medskip
\noindent\textbf{quiz\_pagine\_031\_040\_026}: A. Spiegazione non ancora disponibile.\par\medskip
\noindent\textbf{quiz\_pagine\_041\_050\_020}: A. Risposta ufficiale indicata nella tabella della decima simulazione.\par\medskip
\noindent\textbf{quiz\_pagine\_041\_050\_033}: non verificata / nessuna alternativa corretta. Soluzione non riportata nel PDF; da verificare.\par\medskip
\noindent\textbf{quiz\_pagine\_051\_060\_014}: D. La definizione esprime che $f(x)$ tende a $5$ quando $x$ tende a $+\infty$.\par\medskip
\noindent\textbf{quiz\_pagine\_051\_060\_015}: A. La risposta ufficiale è A; la condizione con un solo $M$ fissato non equivale alla definizione completa di limite.\par\medskip
\noindent\textbf{quiz\_pagine\_061\_070\_003}: A. Il numeratore è equivalente a x²/2 e il denominatore a 3x²; limite 1/6.\par\medskip
\noindent\textbf{quiz\_pagine\_071\_080\_013}: A. Per i limiti notevoli $e^{2x}-1\sim2x$ e $\log(1+3x)\sim3x$, quindi il limite vale $2/3$.\par\medskip
\noindent\textbf{quiz\_pagine\_081\_090\_001}: E. Per $x\to0$, il rapporto è asintotico a $\frac{3x}{25x^2/2}=\frac{6}{25x}$: i limiti laterali hanno segni opposti.\par\medskip
\noindent\textbf{quiz\_pagine\_081\_090\_005}: B. Il denominatore è compreso fra $4$ e $5$ e il numeratore tende a $+\infty$; per confronto il rapporto tende a $+\infty$.\par\medskip
\noindent\textbf{quiz\_pagine\_081\_090\_026}: A. Per $x\to0$, $\log(1+7x)\sim7x$ e $e^{4x}-1\sim4x$, quindi il limite è $7/4$.\par\medskip
\noindent\textbf{quiz\_pagine\_081\_090\_027}: A. $(4^x+2^x)^{1/x}=4(1+2^{-x})^{1/x}\to4$, poiché $\log(1+2^{-x})/x\to0$.\par\medskip
\noindent\textbf{quiz\_pagine\_091\_100\_022}: A. La definizione di limite con $\varepsilon=1$ dà la condizione A, con intorno bucato. Il valore $f(1)$ non è determinato.\par\medskip
\noindent\textbf{quiz\_pagine\_101\_110\_005}: A. Poiché $x^2\to+\infty$ quando $x\to+\infty$, si ha $f(x^2)\to0$. Non è noto il segno.\par\medskip
\noindent\textbf{quiz\_pagine\_101\_110\_009}: A. Usando $2^x-1\sim x\log2$ e $\sin5x\sim5x$, il limite vale $\log2/5$.\par\medskip
\noindent\textbf{quiz\_pagine\_101\_110\_019}: A. Si scrive $x^2-9=(x-3)(x+3)$ e si usa $\sin u/u\to1$: il limite è $x+3\to6$.\par\medskip
\subsubsection*{3.2 ÔÇö Continuit├á di una funzione}
\noindent\textbf{quiz\_pagine\_001\_010\_027}: B. Spiegazione non ancora disponibile.\par\medskip
\noindent\textbf{quiz\_pagine\_011\_020\_022}: B. Per Lagrange esiste $c$ con $f\prime(c)=(0-6)/(6-4)=-3$.\par\medskip
\noindent\textbf{quiz\_pagine\_031\_040\_010}: C. Spiegazione non ancora disponibile.\par\medskip
\noindent\textbf{quiz\_pagine\_041\_050\_013}: E. Risposta ufficiale indicata nella tabella della decima simulazione.\par\medskip
\noindent\textbf{quiz\_pagine\_041\_050\_036}: non verificata / nessuna alternativa corretta. Soluzione non riportata nel PDF; da verificare.\par\medskip
\noindent\textbf{quiz\_pagine\_041\_050\_039}: non verificata / nessuna alternativa corretta. Soluzione non riportata nel PDF; da verificare.\par\medskip
\noindent\textbf{quiz\_pagine\_071\_080\_011}: E. La tabella ufficiale indica E. Nota: anche A è una formulazione vera del teorema degli zeri, ammettendo lo zero agli estremi; il quesito non ha risposta univoca.\par\medskip
\noindent\textbf{quiz\_pagine\_101\_110\_016}: C. La funzione è continua su $[a,10]$ soltanto se $1<a\leq10$. Fra le opzioni la C è l\textbackslash{}textquoteright unica che impone anche un intervallo non degenere. La condizione $a=10$ sarebbe pure ammissibile, ma non è proposta isolatamente.\par\medskip
\subsubsection*{3.4 ÔÇö Confronto tra funzioni: infinitesimi e infiniti}
\noindent\textbf{quiz\_pagine\_001\_010\_012}: D. Spiegazione non ancora disponibile.\par\medskip
\noindent\textbf{quiz\_pagine\_001\_010\_017}: E. Spiegazione non ancora disponibile.\par\medskip
\noindent\textbf{quiz\_pagine\_001\_010\_033}: B. Spiegazione non ancora disponibile.\par\medskip
\noindent\textbf{quiz\_pagine\_001\_010\_039}: C. Spiegazione non ancora disponibile.\par\medskip
\noindent\textbf{quiz\_pagine\_011\_020\_016}: non verificata / nessuna alternativa corretta. Dall\textbackslash{}'ipotesi segue solo $f(x)/x^3\to0$; nessuna delle opzioni qui trascritte lo esprime. La formula dell\textbackslash{}'opzione (c) nel PDF va ricontrollata.\par\medskip
\noindent\textbf{quiz\_pagine\_021\_030\_036}: C. $f(x)/x=[f(x)/x^2]x\to0$.\par\medskip
\noindent\textbf{quiz\_pagine\_031\_040\_009}: D. Spiegazione non ancora disponibile.\par\medskip
\noindent\textbf{quiz\_pagine\_041\_050\_018}: C. Risposta ufficiale indicata nella tabella della decima simulazione.\par\medskip
\noindent\textbf{quiz\_pagine\_051\_060\_010}: C. Lo sviluppo $e^x-\sin x=1+O(x^2)$ mostra che il termine costante deve annullarsi.\par\medskip
\noindent\textbf{quiz\_pagine\_071\_080\_004}: D. Da $f(x)/e^x\to1$ segue $\log f(x)=x+\log(f(x)/e^x)=x+o(1)$, quindi $\log f(x)\sim x$.\par\medskip
\noindent\textbf{quiz\_pagine\_071\_080\_012}: E. Per $x\to0^+$, $\sqrt[4]{x^3+x^4}=x^{3/4}(1+x)^{1/4}\sim x^{3/4}$ e $2x=o(x^{3/4})$.\par\medskip
\noindent\textbf{quiz\_pagine\_071\_080\_019}: non verificata / nessuna alternativa corretta. L'equivalenza $f(x)/x\to1$ per $x\to+\infty$ implica, ponendo $t=1/x$, che $f(1/t)/(1/t)\to1$ per $t\to0^+$; l'affermazione con limite bilatero richiede un'ulteriore ipotesi.\par\medskip
\noindent\textbf{quiz\_pagine\_081\_090\_023}: non verificata / nessuna alternativa corretta. Poiché $f(x)/x=[f(x)/x^2]\,x\to0$, segue $f=o(x)$. Anche l\textbackslash{}textquotesingle alternativa D è vera, perché per $x\to+\infty$ vale $f(1/x)=o(1/x^2)$ e quindi anche $o(x^2)$: pertanto anche D è matematicamente vera.\par\medskip
\noindent\textbf{quiz\_pagine\_091\_100\_008}: A. $\log(1+3x)=3x-\frac92x^2+o(x^2)$, dunque il numeratore è $9x^2+4x^3+o(x^2)$ e il quoziente tende a $0$.\par\medskip
\noindent\textbf{quiz\_pagine\_091\_100\_010}: A. $\sqrt[3]{x^3+x^4}=x(1+x)^{1/3}=x+\frac13x^2+o(x^2)$; quindi $f(x)=\frac23x^2+o(x^2)$.\par\medskip
\noindent\textbf{quiz\_pagine\_091\_100\_013}: C. $\sinh^2(4x)\sim16x^2$ e $1-\cos(3x)\sim9x^2/2$, dunque il limite è $32/9$.\par\medskip
\noindent\textbf{quiz\_pagine\_091\_100\_028}: B. Per $x\to0$, $e^{\sqrt3x}-1\sim\sqrt3x$ e $\sin(\sqrt3x)\sim\sqrt3x$; il rapporto tende a $1$.\par\medskip
\noindent\textbf{quiz\_pagine\_101\_110\_002}: B. Per $x\to0$ la funzione tende a $2$, dunque $f(x)\sim2$.\par\medskip
\subsection*{Capitolo 4}
\subsubsection*{4.1 ÔÇö Derivate}
\noindent\textbf{quiz\_pagine\_001\_010\_001}: E. Spiegazione non ancora disponibile.\par\medskip
\noindent\textbf{quiz\_pagine\_001\_010\_005}: E. Spiegazione non ancora disponibile.\par\medskip
\noindent\textbf{quiz\_pagine\_001\_010\_011}: C. Spiegazione non ancora disponibile.\par\medskip
\noindent\textbf{quiz\_pagine\_011\_020\_008}: D. $f(x)=15\sin^2x$, $f\prime(x)=30\sin x\cos x$, che si annulla in $\pi/2$.\par\medskip
\noindent\textbf{quiz\_pagine\_021\_030\_009}: C. Per $a^{u(x)}$, la derivata \textbackslash{}`e $a^{u(x)}\log(a)u\prime(x)$.\par\medskip
\noindent\textbf{quiz\_pagine\_021\_030\_023}: B. Una primitiva \textbackslash{}`e $\log|f(x)|+C$ dove $f\ne0$; moltiplicare il logaritmo per $2$ ne raddoppia la derivata.\par\medskip
\noindent\textbf{quiz\_pagine\_031\_040\_002}: D. Spiegazione non ancora disponibile.\par\medskip
\noindent\textbf{quiz\_pagine\_031\_040\_005}: E. Spiegazione non ancora disponibile.\par\medskip
\noindent\textbf{quiz\_pagine\_031\_040\_006}: A. Spiegazione non ancora disponibile.\par\medskip
\noindent\textbf{quiz\_pagine\_041\_050\_007}: E. Risposta ufficiale indicata nella tabella della decima simulazione.\par\medskip
\noindent\textbf{quiz\_pagine\_041\_050\_015}: D. Risposta ufficiale indicata nella tabella della decima simulazione.\par\medskip
\noindent\textbf{quiz\_pagine\_041\_050\_032}: non verificata / nessuna alternativa corretta. Soluzione non riportata nel PDF; da verificare.\par\medskip
\noindent\textbf{quiz\_pagine\_051\_060\_016}: D. Il rapporto incrementale in zero è $\sin(1/x)$ e non ha limite.\par\medskip
\noindent\textbf{quiz\_pagine\_061\_070\_002}: D. Imponendo continuità e uguaglianza delle derivate in x=1 si ottiene alfa=4/5, beta=3/5.\par\medskip
\noindent\textbf{quiz\_pagine\_061\_070\_010}: D. Espandendo $e^u=1+u+\frac{u^2}{2}+o(x^2)$, con $u=2x+x^2$, si ottiene $1+2x+3x^2$. Le risposte A e B sono duplicate nel PDF originale.\par\medskip
\noindent\textbf{quiz\_pagine\_061\_070\_014}: A. Si usa log(1+u)=u-u²/2+u³/3+o(x³), con u=x+x².\par\medskip
\noindent\textbf{quiz\_pagine\_061\_070\_022}: E. Sviluppando sqrt(1-x) e poi l’esponenziale si ottiene e(1-x/2-x³/48+o(x³)); soluzione calcolata, tabella ufficiale non presente nell’estratto.\par\medskip
\noindent\textbf{quiz\_pagine\_071\_080\_001}: E. Poiché $g=f^{-2}$, si ha $g'(0)=-2f'(0)/f(0)^3=-6/(-8)=3/4$; la risposta ufficiale è E.\par\medskip
\noindent\textbf{quiz\_pagine\_071\_080\_008}: D. In $0$ entrambi i rami valgono $1$ e le derivate laterali sono entrambe $0$; altrove i rami sono regolari.\par\medskip
\noindent\textbf{quiz\_pagine\_071\_080\_010}: E. Per derivazione composta, $f'(x)=(6x+6x^2)/(3x^2+2x^3)=6(1+x)/(3x+2x^2)$, dove definita.\par\medskip
\noindent\textbf{quiz\_pagine\_071\_080\_016}: C. Per la derivata dell'inversa, se $f'(0)\ne0$ allora $m=1/f'(0)$: con $-1<f'(0)<0$ si ha $m<-1$. Se $f'(0)=0$, la tangente può essere verticale.\par\medskip
\noindent\textbf{quiz\_pagine\_081\_090\_002}: D. $f(e)=e$ e $f\prime(e)=2$, quindi $y=e+2(x-e)=2x-e$.\par\medskip
\noindent\textbf{quiz\_pagine\_081\_090\_008}: E. Per la regola della catena, $(g\circ f)\prime(-2)=g\prime(f(-2))f\prime(-2)=(-7)(-4)=28$. Le alternative A e D sono duplicate nel PDF.\par\medskip
\noindent\textbf{quiz\_pagine\_081\_090\_010}: A. La derivata di $a^x$ è $a^x\log a$ per $a>0$.\par\medskip
\noindent\textbf{quiz\_pagine\_081\_090\_011}: C. $g\prime=f\prime\prime\prime$ esiste, quindi $g$ è derivabile e perciò continua. Non è garantita la continuità di $f\prime\prime\prime$.\par\medskip
\noindent\textbf{quiz\_pagine\_081\_090\_015}: B. Per $x\neq0$, la derivata del primo addendo è $2/(4+x^2)$ e quella del secondo è $-2/(4+x^2)$; sottraendo si ottiene $4/(4+x^2)$.\par\medskip
\noindent\textbf{quiz\_pagine\_081\_090\_021}: D. Per $a=0$, i valori laterali in $0$ sono entrambi $0$ e le derivate laterali sono entrambe $1$; altrove i rami sono regolari.\par\medskip
\noindent\textbf{quiz\_pagine\_081\_090\_024}: non verificata / nessuna alternativa corretta. Per $\sin x<0$, $f\prime(x)=(-\cos x)/(-\sin x)=\cos x/\sin x=1/\tan x$. Le alternative A ed E sono equivalenti, quindi non esiste una sola risposta corretta.\par\medskip
\noindent\textbf{quiz\_pagine\_091\_100\_003}: E. La regola della catena dà $f\prime(x)=\frac13(e^{4x}\cos^8x)^{-2/3}e^{4x}(4\cos^8x-8\cos^7x\sin x)$, equivalente alla E, con prolungamento nei punti $\cos x=0$.\par\medskip
\noindent\textbf{quiz\_pagine\_091\_100\_016}: non verificata / nessuna alternativa corretta. Le derivate di ordine dispari di $\cos x$ si annullano in $0$: A è vera; anche D è vera perché descrive un sottoinsieme degli ordini dispari. Il testo ufficiale indica A, ma il quesito ha più risposte vere.\par\medskip
\noindent\textbf{quiz\_pagine\_091\_100\_017}: E. $f\prime(x)=-5e^{-5x}/(3-e^{-5x})^2=-5e^{5x}/(3e^{5x}-1)^2$.\par\medskip
\noindent\textbf{quiz\_pagine\_091\_100\_023}: A. Il teorema di prolungamento della derivata: continuità in $2$ e limite finito di $f\prime$ implicano derivabilità in $2$ con $f\prime(2)=4$.\par\medskip
\noindent\textbf{quiz\_pagine\_101\_110\_006}: E. Con la regola del quoziente, $f\prime(x)=\frac{-e^{1/x}/x- (e^{1/x}-1)}{x^2}=\frac{x-(x+1)e^{1/x}}{x^3}$.\par\medskip
\noindent\textbf{quiz\_pagine\_101\_110\_012}: E. Le derivate di ordine dispari del coseno sono proporzionali al seno, e quindi si annullano in $0$. Si intende $f^n$ come derivata di ordine $n$.\par\medskip
\noindent\textbf{quiz\_pagine\_101\_110\_014}: A. La continuità in $0$ richiede $a=\cos0=1$. Le derivate laterali sono entrambe nulle, dunque la funzione è derivabile anche in $0$.\par\medskip
\subsubsection*{4.2 ÔÇö I teoremi fondamentali del calcolo differenziale}
\noindent\textbf{quiz\_pagine\_001\_010\_038}: A. Spiegazione non ancora disponibile.\par\medskip
\noindent\textbf{quiz\_pagine\_011\_020\_012}: non verificata / nessuna alternativa corretta. $f(2)=2>0$, $f\prime(2)=-3<0$, $f\prime\prime(2)=-10<0$: positiva, decrescente e concava. Nessuna opzione coincide.\par\medskip
\noindent\textbf{quiz\_pagine\_011\_020\_015}: non verificata / nessuna alternativa corretta. La funzione \textbackslash{}`e $e+O(x^4)$, quindi il polinomio di ordine 3 \textbackslash{}`e costante $e$. Verificare la formula stampata nel PDF: l\textbackslash{}'estrazione non chiarisce l\textbackslash{}'esponente.\par\medskip
\noindent\textbf{quiz\_pagine\_011\_020\_021}: C. $1/\cos x=1+x^2/2+o(x^2)$ e $\sin x=x+o(x^2)$.\par\medskip
\noindent\textbf{quiz\_pagine\_011\_020\_023}: D. $3x^2(1-3x+9x^2)-3x^2-3x^4=-9x^3+24x^4+o(x^4)$.\par\medskip
\noindent\textbf{quiz\_pagine\_011\_020\_025}: A. $f(x)=2+\log5-x^2/5+o(x^2)$: massimo locale stretto.\par\medskip
\noindent\textbf{quiz\_pagine\_011\_020\_026}: E. $\sinh(\log(1+x))=((1+x)-(1+x)^{-1})/2=x-x^2/2+x^3/2+o(x^3)$.\par\medskip
\noindent\textbf{quiz\_pagine\_011\_020\_029}: C. $f(7)=-3$ e $f\prime(7)=2$, quindi la tangente \textbackslash{}`e $y=-3+2(x-7)$.\par\medskip
\noindent\textbf{quiz\_pagine\_011\_020\_030}: A. $\cosh x=1+x^2/2+O(x^4)$, dunque $\log\cosh x=x^2/2+O(x^4)$.\par\medskip
\noindent\textbf{quiz\_pagine\_011\_020\_032}: A. $\sinh x=x+x^3/6+o(x^3)$; in $\log(1+u)=u-u^2/2+u^3/3+\cdots$ il coefficiente cubico \textbackslash{}`e $1/6+1/3=1/2$.\par\medskip
\noindent\textbf{quiz\_pagine\_011\_020\_033}: D. $2e^{-x}-x^2=2-2x+O(x^3)$ e $(1+x^2)^{-1}=1-x^2+o(x^2)$.\par\medskip
\noindent\textbf{quiz\_pagine\_011\_020\_034}: E. Il termine dominante $-4(x-2)^6$ d\textbackslash{}`a un massimo piatto in $2$, non un flesso crescente.\par\medskip
\noindent\textbf{quiz\_pagine\_011\_020\_035}: C. $(1+2x^2)(1+x^2+x^4+\cdots)=1+3x^2+3x^4+\cdots$.\par\medskip
\noindent\textbf{quiz\_pagine\_021\_030\_001}: B. $f(1)=e^{-1/3}$, $f\prime(1)=-2e^{-1/3}/3$, $f\prime\prime(1)=-2e^{-1/3}/9$, dunque il coefficiente quadratico \textbackslash{}`e $-e^{-1/3}/9$.\par\medskip
\noindent\textbf{quiz\_pagine\_021\_030\_004}: C. $\cosh x=1+x^2/2+x^4/24+o(x^4)$; $e^{\cosh x}=e[1+x^2/2+(1/24+1/8)x^4]+o(x^4)$.\par\medskip
\noindent\textbf{quiz\_pagine\_021\_030\_012}: D. $3^{2\sin x}=e^{2\log3(x+o(x^2))}=1+2x\log3+2x^2\log^23+o(x^2)$.\par\medskip
\noindent\textbf{quiz\_pagine\_021\_030\_028}: B. $(1+x^4)^{1/3}=1+x^4/3+o(x^4)$, $(1+x^3)^{-1/4}=1-x^3/4+o(x^4)$.\par\medskip
\noindent\textbf{quiz\_pagine\_041\_050\_002}: E. Risposta ufficiale indicata nella tabella della decima simulazione.\par\medskip
\noindent\textbf{quiz\_pagine\_041\_050\_003}: D. Risposta ufficiale indicata nella tabella della decima simulazione.\par\medskip
\noindent\textbf{quiz\_pagine\_041\_050\_006}: D. Per il teorema di Rolle, poiché f(-3)=f(-2)=0, esiste c interno con derivata nulla.\par\medskip
\noindent\textbf{quiz\_pagine\_041\_050\_017}: E. Risposta ufficiale indicata nella tabella della decima simulazione.\par\medskip
\noindent\textbf{quiz\_pagine\_051\_060\_005}: E. Il teorema di Rolle richiede continuità sul chiuso, derivabilità nell’aperto e uguaglianza dei valori agli estremi.\par\medskip
\noindent\textbf{quiz\_pagine\_051\_060\_006}: B. La risposta ufficiale indica B; le ipotesi riportate non specificano una forma indeterminata per applicare de l’Hôpital.\par\medskip
\noindent\textbf{quiz\_pagine\_051\_060\_012}: E. La risposta ufficiale è E, ma A e B risultano soddisfatte: la funzione è continua sul chiuso e derivabile nell’aperto.\par\medskip
\noindent\textbf{quiz\_pagine\_061\_070\_023}: D. Per Rolle, f(-2)=f(0)=0 implica un punto c in (-2,0) con derivata nulla.\par\medskip
\noindent\textbf{quiz\_pagine\_071\_080\_003}: E. Nello sviluppo in $x=3$ il coefficiente di $(x-3)^2$ è $3$, dunque $f''(3)=6$. La tabella ufficiale indica E per la derivata terza; la sua esistenza richiede cautela con il solo resto di Peano.\par\medskip
\noindent\textbf{quiz\_pagine\_071\_080\_014}: C. Lo sviluppo centrato in $-4$ dà $f'(-4)=0$ e $f''(-4)=6>0$: $-4$ è un minimo relativo.\par\medskip
\noindent\textbf{quiz\_pagine\_081\_090\_003}: B. $\cos x=1-x^2/2+O(x^4)$ e $\cosh x=1+x^2/2+O(x^4)$: il termine quadratico si annulla, quindi $f(x)=1+o(x^2)$. La risposta C è uno sviluppo corretto di ordine superiore.\par\medskip
\noindent\textbf{quiz\_pagine\_081\_090\_019}: E. $e^x-1=x+O(x^2)$, quindi $(e^x-1)^2=x^2+O(x^3)$ e $\sin x=x+O(x^3)$; il polinomio è $x+x^2$. L\textbackslash{}textquotesingle alternativa A rappresenta invece lo sviluppo con resto.\par\medskip
\noindent\textbf{quiz\_pagine\_091\_100\_002}: B. Il termine principale è $2(x+3)^3$, che cambia segno attraversando $-3$. Non è lecito dedurre $f(0)$ dallo sviluppo locale; $f^{\prime\prime\prime}(-3)=12$ se la derivata esiste.\par\medskip
\noindent\textbf{quiz\_pagine\_091\_100\_012}: D. Dallo sviluppo $f(\pi)=\pi$ e $f\prime(\pi)=2$; la tangente è $y=\pi+2(x-\pi)$.\par\medskip
\noindent\textbf{quiz\_pagine\_091\_100\_015}: C. Per il teorema di Rolle esiste almeno un $c\in(0,3)$ con $f\prime(c)=0$; il numero esatto di zeri non è determinato.\par\medskip
\noindent\textbf{quiz\_pagine\_091\_100\_019}: non verificata / nessuna alternativa corretta. Poiché $\cos x=1-x^2/2+O(x^4)$ e $\cosh x=1+x^2/2+O(x^4)$, il prodotto è $1+o(x^2)$. Anche B è formalmente vera: $x^4=o(x^2)$, quindi la scelta del coefficiente di $x^4$ non cambia il resto $o(x^2)$. La tabella ufficiale indica E.\par\medskip
\noindent\textbf{quiz\_pagine\_101\_110\_010}: C. Da $\sin x=x-x^3/6+o(x^3)$ e $x^3\cos x=x^3+o(x^3)$ si ottiene $1+x-7x^3/6+o(x^3)$.\par\medskip
\noindent\textbf{quiz\_pagine\_101\_110\_020}: E. Ponendo $u=\sin x/3=x/3-x^3/18+o(x^3)$ e usando $\log(1+u)=u-u^2/2+u^3/3+o(u^3)$, si trova $x/3-x^2/18-7x^3/162+o(x^3)$.\par\medskip
\subsubsection*{4.3 ÔÇö Propriet├á delle funzioni: monotonia e punti estremali}
\noindent\textbf{quiz\_pagine\_001\_010\_013}: A. Spiegazione non ancora disponibile.\par\medskip
\noindent\textbf{quiz\_pagine\_001\_010\_014}: B. Spiegazione non ancora disponibile.\par\medskip
\noindent\textbf{quiz\_pagine\_001\_010\_024}: A. Spiegazione non ancora disponibile.\par\medskip
\noindent\textbf{quiz\_pagine\_001\_010\_032}: E. Spiegazione non ancora disponibile.\par\medskip
\noindent\textbf{quiz\_pagine\_011\_020\_024}: B. $(f^{-1})\prime(f(2))=1/f\prime(2)=7$.\par\medskip
\noindent\textbf{quiz\_pagine\_021\_030\_003}: D. $f\prime=3(x-1)(x-3)$: nell\textbackslash{}'intervallo il punto critico \textbackslash{}`e $1$. $f(-2)=-49$, $f(1)=5$, $f(2)=3$.\par\medskip
\noindent\textbf{quiz\_pagine\_031\_040\_004}: D. Spiegazione non ancora disponibile.\par\medskip
\noindent\textbf{quiz\_pagine\_041\_050\_009}: E. Risposta ufficiale indicata nella tabella della decima simulazione.\par\medskip
\noindent\textbf{quiz\_pagine\_041\_050\_019}: C. Risposta ufficiale indicata nella tabella della decima simulazione.\par\medskip
\noindent\textbf{quiz\_pagine\_041\_050\_043}: non verificata / nessuna alternativa corretta. Soluzione non riportata nel PDF; da verificare.\par\medskip
\noindent\textbf{quiz\_pagine\_051\_060\_011}: D. Il termine $1/(1+x^2)$ raggiunge il massimo $1$ in $x=0$.\par\medskip
\noindent\textbf{quiz\_pagine\_051\_060\_013}: C. Poiché $1/(1+x^2)>0$ e tende a zero per $|x|	o\infty$, l’estremo inferiore è $\log2$.\par\medskip
\noindent\textbf{quiz\_pagine\_061\_070\_024}: non verificata / nessuna alternativa corretta. Classificazione basata su estremi relativi; soluzione non riportata nelle pagine fornite. Notazione dei limiti con segni accoppiati da verificare.\par\medskip
\noindent\textbf{quiz\_pagine\_081\_090\_020}: B. $f\prime(x)=3^x\log3+3x^2>0$, dunque $f$ è strettamente crescente. Per continuità e perché l\textbackslash{}textquotesingle immagine è $\mathbb R$, il valore indicato è assunto esattamente una volta.\par\medskip
\noindent\textbf{quiz\_pagine\_091\_100\_007}: A. $f\prime(x)=5x^4+1>0$, quindi ha un unico zero; poiché $f(-2)<0<f(2)$, lo zero è in $(-2,2)$.\par\medskip
\noindent\textbf{quiz\_pagine\_091\_100\_011}: D. $f\prime(x)=7x^6+7e^{7x}>0$ e $f$ è suriettiva su $\mathbb R$. Poiché $\bar y=f(5)+2>f(5)$, la soluzione unica soddisfa $x_0>5$.\par\medskip
\noindent\textbf{quiz\_pagine\_101\_110\_017}: non verificata / nessuna alternativa corretta. La tabella indica A e i coefficienti suggeriscono $f(2)=2$, $f\prime(2)=6$, $f\prime\prime(2)=2$. Tuttavia il solo resto $o((x-2)^2)$ non garantisce convessità né monotonia in un intero intorno senza ulteriori ipotesi di regolarità: soluzione rigorosamente non determinabile.\par\medskip
\subsubsection*{4.4 ÔÇö Lo studio di funzione}
\noindent\textbf{quiz\_pagine\_001\_010\_007}: C. Spiegazione non ancora disponibile.\par\medskip
\noindent\textbf{quiz\_pagine\_001\_010\_023}: D. Spiegazione non ancora disponibile.\par\medskip
\noindent\textbf{quiz\_pagine\_011\_020\_006}: B. $\log(3+2e^x)=x+\log(2+3e^{-x})=x+\log2+o(1)$.\par\medskip
\noindent\textbf{quiz\_pagine\_041\_050\_011}: B. La formula centrale è stata ricostruita dal testo estratto: verificare visivamente sul PDF.\par\medskip
\noindent\textbf{quiz\_pagine\_061\_070\_006}: B. f(x)-x=log x diverge, quindi non esiste intercetta finita per un asintoto obliquo.\par\medskip
\noindent\textbf{quiz\_pagine\_071\_080\_006}: B. Per $x\to+\infty$, $\log(e^{3x}+4)=3x+o(1)$ e $(\cos x-1)/x=o(1)$; l'asintoto è $y=7x$.\par\medskip
\noindent\textbf{quiz\_pagine\_091\_100\_006}: non verificata / nessuna alternativa corretta. La tabella ufficiale indica E. La funzione è non negativa in $(0,5)$ e $x=2$ è uno zero di ordine pari, dunque è un minimo locale, non un massimo: anche A risulta falsa. Quesito originale problematico, con almeno due risposte NON soddisfatte.\par\medskip
\noindent\textbf{quiz\_pagine\_091\_100\_024}: B. Per $x\to-\infty$, $e^x\to0$ mentre $\cos x$ oscilla fra $-1$ e $1$. Su infiniti intervalli negativi cambia segno, quindi la funzione ha infiniti zeri.\par\medskip
\subsection*{Capitolo 5}
\subsubsection*{5.2 ÔÇö Propriet├á dellÔÇÖintegrale}
\noindent\textbf{quiz\_pagine\_011\_020\_018}: A. $F\prime(x)=x^4\cos^4x\ge0$, quindi $F$ \textbackslash{}`e non decrescente su $\mathbb R$.\par\medskip
\noindent\textbf{quiz\_pagine\_031\_040\_007}: B. Spiegazione non ancora disponibile.\par\medskip
\noindent\textbf{quiz\_pagine\_041\_050\_010}: A. Risposta ufficiale indicata nella tabella della decima simulazione.\par\medskip
\noindent\textbf{quiz\_pagine\_041\_050\_038}: non verificata / nessuna alternativa corretta. Soluzione non riportata nel PDF; da verificare.\par\medskip
\noindent\textbf{quiz\_pagine\_041\_050\_040}: non verificata / nessuna alternativa corretta. Soluzione non riportata nel PDF; da verificare.\par\medskip
\noindent\textbf{quiz\_pagine\_051\_060\_002}: E. Una potenza dispari di una funzione dispari resta dispari e ha integrale nullo su un intervallo simmetrico.\par\medskip
\noindent\textbf{quiz\_pagine\_051\_060\_003}: D. Il teorema fondamentale del calcolo richiede la continuità di $f$ per garantire $F\prime=f$ nell’interno.\par\medskip
\noindent\textbf{quiz\_pagine\_051\_060\_009}: C. Per $F(x)=\int_x^a f(t)dt$ vale $F\prime(x)=-f(x)\le0$; inoltre $F(a)=0$.\par\medskip
\noindent\textbf{quiz\_pagine\_051\_060\_019}: C. Poiché $F\prime(x)=-f(x)\le0$, la funzione integrale è non crescente.\par\medskip
\noindent\textbf{quiz\_pagine\_061\_070\_008}: E. DISCREPANZA: risposta ufficiale E, ma integrando una funzione pari le opzioni A e D sono vere; E è falsa in generale (esempio f=1).\par\medskip
\noindent\textbf{quiz\_pagine\_061\_070\_011}: B. F è strettamente crescente e F(0)=0, perciò F(-1)<0.\par\medskip
\noindent\textbf{quiz\_pagine\_061\_070\_018}: A. Integrando le disuguaglianze strette 2<f<4 e dividendo per b-a>0 si ottiene 2<m<4.\par\medskip
\noindent\textbf{quiz\_pagine\_071\_080\_005}: A. Per $x\ge x_0$ l'integrale di una funzione non negativa è non negativo: $F(x)=\int_{x_0}^x f(t)\,dt\ge0$.\par\medskip
\noindent\textbf{quiz\_pagine\_091\_100\_021}: A. $F(x)=\log\cosh x\geq0$ per ogni $x$, e $F(x)>0$ anche per $x<0$; dunque A è falsa. $F\prime=\tanh x$ e $F^{\prime\prime}=1/\cosh^2x>0$.\par\medskip
\noindent\textbf{quiz\_pagine\_101\_110\_008}: D. Per il teorema fondamentale del calcolo $F\prime(x)=e^{-x^2}>0$: $F$ è strettamente crescente. La funzione integranda è pari, perciò $F$ è dispari.\par\medskip
\subsubsection*{5.3 ÔÇö Integrali indefiniti e regole di integrazione}
\noindent\textbf{quiz\_pagine\_041\_050\_034}: non verificata / nessuna alternativa corretta. Soluzione non riportata nel PDF; da verificare.\par\medskip
\noindent\textbf{quiz\_pagine\_041\_050\_041}: non verificata / nessuna alternativa corretta. Soluzione non riportata nel PDF; da verificare.\par\medskip
\noindent\textbf{quiz\_pagine\_081\_090\_004}: A. Una primitiva è $-\frac32\cos 2x$, assente fra le alternative B--E.\par\medskip
\noindent\textbf{quiz\_pagine\_081\_090\_009}: D. Una primitiva è $-\cos(f(x))$; agli estremi si ottiene $\cos f(0)-\cos f(2)$.\par\medskip
\subsubsection*{5.4 ÔÇö Tecniche di integrazione}
\noindent\textbf{quiz\_pagine\_001\_010\_034}: C. Spiegazione non ancora disponibile.\par\medskip
\noindent\textbf{quiz\_pagine\_001\_010\_037}: A. Spiegazione non ancora disponibile.\par\medskip
\noindent\textbf{quiz\_pagine\_011\_020\_017}: C. Per simmetria, $4\int_0^{\pi/2}\cos x\,dx=4$.\par\medskip
\noindent\textbf{quiz\_pagine\_021\_030\_008}: A. $\frac{x-4}{5-x}=-1-\frac1{x-5}$ e la derivata \textbackslash{}`e $1/(x-5)^2$.\par\medskip
\noindent\textbf{quiz\_pagine\_021\_030\_022}: E. La (e) non \textbackslash{}`e definita su tutta $\mathbb R$ e la sua derivata include $-1/(3x)$.\par\medskip
\noindent\textbf{quiz\_pagine\_021\_030\_025}: D. Il testo usa la definizione di primitiva su un intervallo: $\mathbb R\setminus\{0\}$ non \textbackslash{}`e un intervallo, pur essendo la derivata $1/x$ su ciascuna componente.\par\medskip
\noindent\textbf{quiz\_pagine\_041\_050\_035}: non verificata / nessuna alternativa corretta. Soluzione non riportata nel PDF; da verificare.\par\medskip
\noindent\textbf{quiz\_pagine\_061\_070\_005}: E. Integrazione per parti con u=x e dv=g\textbackslash{}prime(x)dx: 2g(2)-integrale di g.\par\medskip
\noindent\textbf{quiz\_pagine\_061\_070\_007}: C. Sostituendo u=3t+2 si ha dt=du/3, con estremi 2 e 8.\par\medskip
\noindent\textbf{quiz\_pagine\_061\_070\_012}: E. Con t=2x, dx=dt/2 e gli estremi diventano 0 e 2.\par\medskip
\noindent\textbf{quiz\_pagine\_061\_070\_017}: A. Per parti con dv=dx/x. DISCREPANZA: A e B sono equivalenti su [1,2]; la tabella ufficiale indica A.\par\medskip
\noindent\textbf{quiz\_pagine\_061\_070\_019}: D. Con t=log x e dt=dx/x gli estremi diventano 0 e log 2. DISCREPANZA: anche B ripete identicamente il membro sinistro.\par\medskip
\subsubsection*{5.5 ÔÇö Integrali impropri}
\noindent\textbf{quiz\_pagine\_001\_010\_029}: A. Spiegazione non ancora disponibile.\par\medskip
\noindent\textbf{quiz\_pagine\_011\_020\_019}: A. Vicino a zero, $1-\cos x\sim x^2/2$ e l\textbackslash{}'integranda \textbackslash{}`e equivalente a $2/x$.\par\medskip
\noindent\textbf{quiz\_pagine\_041\_050\_008}: D. Risposta ufficiale indicata nella tabella della decima simulazione.\par\medskip
\noindent\textbf{quiz\_pagine\_051\_060\_007}: D. Per il criterio del confronto, se $0\le f\le g$ e l’integrale di $g$ converge, converge anche quello di $f$.\par\medskip
\noindent\textbf{quiz\_pagine\_051\_060\_017}: A. La presenza della singolarità $x=2$ impedisce la convergenza per intervalli che la contengono.\par\medskip
\noindent\textbf{quiz\_pagine\_051\_060\_018}: C. Vicino a zero l’integranda si comporta come $x^{1-\alpha}$, all’infinito come una costante per $x^\alpha$.\par\medskip
\noindent\textbf{quiz\_pagine\_051\_060\_020}: A. Vicino a zero $\sin x/(x\sqrt{x})\sim1/\sqrt{x}$; all’infinito converge per confronto assoluto.\par\medskip
\noindent\textbf{quiz\_pagine\_061\_070\_009}: A. Per alfa positivo, vicino a zero serve alfa<1 e all’infinito serve 2alfa>1.\par\medskip
\noindent\textbf{quiz\_pagine\_061\_070\_013}: D. Una primitiva è -3x\textasciicircum{}(-1/3); il limite superiore vale zero e resta 3/a\textasciicircum{}(1/3).\par\medskip
\noindent\textbf{quiz\_pagine\_061\_070\_015}: C. Vicino a zero sin t/t\textasciicircum{}alfa è equivalente a t\textasciicircum{}(1-alfa); diverge se alfa>=2.\par\medskip
\noindent\textbf{quiz\_pagine\_061\_070\_016}: A. La primitiva è arctan t e il salto fra gli estremi infiniti vale pi.\par\medskip
\noindent\textbf{quiz\_pagine\_061\_070\_020}: B. Poiché arctan x tende a pi/2, il comportamento è quello di x\textasciicircum{}alfa; convergenza per alfa<-1.\par\medskip
\noindent\textbf{quiz\_pagine\_071\_080\_002}: B. L'integranda è integrabile in $0$ e per $x\to+\infty$ è equivalente a $x^{-3/2}$; l'integrale converge ed è positivo, come indicato nella soluzione ufficiale.\par\medskip
\noindent\textbf{quiz\_pagine\_071\_080\_015}: A. Per $x\to+\infty$ il termine $1/\sqrt{x}$ ha integrale divergente e l'altro addendo è non negativo; diverge l'integrale totale.\par\medskip
\noindent\textbf{quiz\_pagine\_081\_090\_007}: D. Per $x\to0^+$, $1/(\sqrt{x}+x^2)\sim x^{-1/2}$, integrabile in $0$; l\textbackslash{}textquotesingle integranda è positiva, quindi il valore è positivo.\par\medskip
\noindent\textbf{quiz\_pagine\_081\_090\_018}: A. Per il criterio del confronto, $0\leq f\leq g$ e la convergenza di $\int g$ implicano la convergenza di $\int f$. La formulazione D è ambigua: un integrale improprio di una funzione con segno variabile può non convergere, ma il termine ``indeterminato non è definito.\par\medskip
\noindent\textbf{quiz\_pagine\_081\_090\_025}: C. L\textbackslash{}textquotesingle integranda è $(1-\pi)x^{3-\pi}$: poiché $3-\pi>-1$, l\textbackslash{}textquotesingle integrale converge e vale $(1-\pi)/(4-\pi)<0$.\par\medskip
\noindent\textbf{quiz\_pagine\_091\_100\_020}: B. Una primitiva è $3/[2(x-3)^2]$. Il valore $F(4)-F(3^+)$ diverge a $-\infty$.\par\medskip
\noindent\textbf{quiz\_pagine\_101\_110\_003}: C. Una primitiva è $\frac12(\log x)^2$, che tende a $+\infty$ per $x\to+\infty$.\par\medskip
\subsection*{Capitolo 6}
\subsubsection*{6.1 ÔÇö Metodi risolutivi di alcune equazioni differenziali ordinarie}
\noindent\textbf{quiz\_pagine\_001\_010\_008}: D. Spiegazione non ancora disponibile.\par\medskip
\noindent\textbf{quiz\_pagine\_011\_020\_020}: D. Il polinomio caratteristico \textbackslash{}`e $(r-2)^2$, radice doppia $r=2$.\par\medskip
\noindent\textbf{quiz\_pagine\_021\_030\_021}: D. $y\prime=-\sin x$, quindi $\sin^2x+\cos^2x=1$.\par\medskip
\noindent\textbf{quiz\_pagine\_031\_040\_008}: A. Spiegazione non ancora disponibile.\par\medskip
\noindent\textbf{quiz\_pagine\_041\_050\_037}: non verificata / nessuna alternativa corretta. Soluzione non riportata nel PDF; da verificare.\par\medskip
\noindent\textbf{quiz\_pagine\_041\_050\_042}: non verificata / nessuna alternativa corretta. Soluzione non riportata nel PDF; da verificare.\par\medskip
\noindent\textbf{quiz\_pagine\_071\_080\_009}: B. L'equazione $x'=x-6$ ha soluzione $x(t)=6+Ce^t$; da $x(-1)=5$ segue $C=-e$, dunque $x(t)=6-e^{t+1}$.\par\medskip
\noindent\textbf{quiz\_pagine\_071\_080\_018}: A. L'equazione $y''-(2/3)y'=0$ ha soluzione $y=C_1+C_2e^{2x/3}$. Il vincolo $y(0)=0$ impone $C_1=-C_2$, lasciando infiniti valori di $C_2$.\par\medskip
\noindent\textbf{quiz\_pagine\_081\_090\_012}: C. Per $y=x$ e $x\neq0,1$ il secondo membro è $1$, uguale a $y\prime=1$; nei punti $y=0,1$ il denominatore si annulla.\par\medskip
\noindent\textbf{quiz\_pagine\_081\_090\_022}: C. Si scrive $y\prime-x^2y=0$ (lineare) e, per $y\neq0$, $dy/y=x^2dx$ (separabile). La soluzione nulla è anch\textbackslash{}textquotesingle essa ammessa.\par\medskip
\noindent\textbf{quiz\_pagine\_091\_100\_014}: D. Per linearità e omogeneità dell'equazione, ogni combinazione lineare di soluzioni è soluzione. Una costante additiva non lo è necessariamente.\par\medskip
\noindent\textbf{quiz\_pagine\_091\_100\_027}: C. Le soluzioni sono $y=A\cos(x/2)+B\sin(x/2)$. La sola soluzione costante è $y\equiv0$; imporre un unico valore iniziale non determina una soluzione unica.\par\medskip
\noindent\textbf{quiz\_pagine\_101\_110\_007}: B. Sostituendo $y=e^2\sin t$ si trova $-e^2\sin t=ae^2\sin t+b$, dunque $a=-1$, $b=0$.\par\medskip
\noindent\textbf{quiz\_pagine\_101\_110\_011}: D. L\textbackslash{}textquoteright equazione caratteristica è $\lambda(\lambda+1/5)=0$, perciò $y=C_1+C_2e^{-t/5}$, non una combinazione di seni e coseni.\par\medskip
\subsubsection*{6.2 ÔÇö Buona positura del problema di Cauchy}
\noindent\textbf{quiz\_pagine\_041\_050\_014}: A. Risposta ufficiale indicata nella tabella della decima simulazione.\par\medskip
\noindent\textbf{quiz\_pagine\_051\_060\_001}: A. Per $a=1$ il dato iniziale non appartiene al dominio reale $x\ge 2$ della radice.\par\medskip
\subsection*{Capitolo 7}
\subsubsection*{7.1 ÔÇö Concetti di base}
\noindent\textbf{quiz\_pagine\_091\_100\_009}: C. Le radici quadrate di $i$ sono opposte: $z_1=e^{i\pi/4}$ e $z_2=-z_1$, quindi $z_1/z_2=-1\in\mathbb R$.\par\medskip
\noindent\textbf{quiz\_pagine\_091\_100\_025}: B. Poiché i coefficienti sono reali, anche $2-i$ è radice. Le tre radici distinte $0,2+i,2-i$ impongono grado almeno $3$.\par\medskip
\end{document}
